\documentclass[../DoAn.tex]{subfiles}
\begin{document}

Chương này trình bày khảo sát hiện trạng và phân tích yêu cầu cho trò chơi giải đố từ vựng dạng ô chữ. Nội dung chương tập trung làm rõ nhu cầu xây dựng sản phẩm, phạm vi hệ thống, các chức năng chính cần có và các ràng buộc phi chức năng ảnh hưởng trực tiếp tới thiết kế.

\section{Khảo sát hiện trạng}
\label{section:2.1}
Các trò chơi dạng ô chữ phổ biến thường được thiết kế xoay quanh tiếng Anh hoặc các ngôn ngữ không dấu. Với hướng tiếp cận này, mỗi ô trên bảng thường tương ứng trực tiếp với một chữ cái, người chơi nhập từng ký tự hoặc nhập toàn bộ đáp án, sau đó hệ thống so sánh với từ đích. Khi áp dụng cho tiếng Việt, cách tiếp cận này gặp khó khăn do một chữ cái có thể mang dấu thanh hoặc biến thể nguyên âm như ``ă'', ``â'', ``ê'', ``ô'', ``ơ'', ``ư''. Ngoài ra, nhiều đáp án tiếng Việt là cụm từ gồm nhiều tiếng, ví dụ có khoảng trắng trong cách viết thông thường nhưng lại cần được xếp liền nhau trên bảng ô chữ.

Trong phạm vi phân tích yêu cầu của đồ án, có thể xem xét một số hướng tiếp cận phổ biến như ô chữ truyền thống, flashcard, câu hỏi trắc nghiệm hoặc ghép chữ đơn giản. Các hướng tiếp cận này đều hỗ trợ luyện từ ở một mức độ nhất định, nhưng chưa đồng thời giải quyết bài toán tổ chức từ vựng theo chủ đề, sinh bảng ô chữ tự động, kiểm tra đáp án tiếng Việt đầy đủ và lưu lại tiến độ của người chơi.

Đối với người học hoặc người chơi phổ thông, một trò chơi giải đố từ vựng dạng ô chữ cần đáp ứng đồng thời hai nhóm yêu cầu. Nhóm thứ nhất là yêu cầu nội dung, bao gồm từ vựng đủ đa dạng, gợi ý dễ hiểu, chủ đề gần gũi và độ dài từ phù hợp với kích thước bảng. Nhóm thứ hai là yêu cầu tương tác, bao gồm thao tác chọn gợi ý nhanh, nhập đáp án thuận tiện, phản hồi đúng sai rõ ràng và khả năng quay lại chơi tiếp. Nếu thiếu một trong hai nhóm yêu cầu này, trò chơi dễ trở thành một bảng ô chữ tĩnh hoặc một danh sách câu hỏi rời rạc, không tạo được trải nghiệm chơi liên tục.

\begin{table}[H]
\centering
\begin{tabular}{|p{3.2cm}|p{5.2cm}|p{5.2cm}|}
\hline
\textbf{Hướng tiếp cận} & \textbf{Ưu điểm} & \textbf{Hạn chế đối với đề tài} \\ \hline
Ô chữ truyền thống trên giấy & Dễ hiểu, không cần thiết bị, phù hợp với hoạt động học tập cơ bản & Không tự động kiểm tra đáp án, không lưu tiến độ, khó cá nhân hóa nội dung và khó sinh nhiều màn chơi \\ \hline
Ô chữ tiếng Anh & Có nhiều cơ chế chơi hoàn thiện, giao diện quen thuộc, luật chơi rõ ràng & Không xử lý trực tiếp đặc thù dấu tiếng Việt và cụm từ có khoảng trắng \\ \hline
Flashcard từ vựng & Phù hợp học nghĩa của từ, dễ tổ chức theo chủ đề & Ít yếu tố giao nhau giữa các từ, không tạo được bài toán suy luận dạng ô chữ \\ \hline
Ghép chữ đơn giản & Tương tác nhanh, dễ chơi trên thiết bị cá nhân & Thường chỉ xử lý từ đơn hoặc ký tự rời, chưa nhấn mạnh gợi ý ngữ nghĩa và tiến độ theo màn \\ \hline
\end{tabular}
\caption{Phân tích một số hướng tiếp cận liên quan đến trò chơi giải đố từ vựng}
\label{table:related_products}
\end{table}

Từ bảng so sánh trên, đồ án lựa chọn hướng xây dựng một trò chơi giải đố từ vựng dạng ô chữ chuyên biệt cho tiếng Việt thay vì chỉ mô phỏng lại trò chơi ô chữ tiếng Anh. Trọng tâm của hệ thống là xử lý từ tiếng Việt đúng cách, tổ chức dữ liệu theo chủ đề và tạo trải nghiệm chơi có tiến độ. Do phạm vi đồ án là một sản phẩm ứng dụng, các yêu cầu được ưu tiên theo mức độ ảnh hưởng tới khả năng chơi được của nguyên mẫu.

\section{Tổng quan chức năng}
\label{section:2.2}
Tác nhân chính của hệ thống là người chơi. Người chơi có thể chọn hoặc tạo tài khoản cục bộ, chọn chủ đề, chọn màn chơi, xem danh sách gợi ý, nhập đáp án, sử dụng sao để nhận gợi ý, xem điểm số và theo dõi tiến độ hoàn thành. Bên trong phạm vi hệ thống, ứng dụng tự động quản lý dữ liệu chủ đề, sinh bảng ô chữ, kiểm tra đáp án, tính điểm, cập nhật nhiệm vụ và lưu dữ liệu tiến độ. Các hoạt động nội bộ này không được xem là tác nhân ngoài mà là trách nhiệm xử lý của phần mềm.

\begin{table}[H]
\centering
\begin{tabular}{|p{3.2cm}|p{10.8cm}|}
\hline
\textbf{Thành phần} & \textbf{Mô tả vai trò} \\ \hline
Người chơi & Tác nhân ngoài trực tiếp sử dụng ứng dụng để chọn tài khoản, chọn chủ đề, giải ô chữ, dùng gợi ý, xem kết quả và quay lại các màn đã chơi. \\ \hline
Phạm vi hệ thống & Bao gồm các chức năng nội bộ: sinh bảng ô chữ, hiển thị gợi ý, kiểm tra đáp án, tính điểm, quản lý sao, cập nhật nhiệm vụ, lưu tiến độ và khôi phục dữ liệu khi người chơi mở lại ứng dụng. \\ \hline
\end{tabular}
\caption{Tác nhân và phạm vi hệ thống}
\label{table:actors}
\end{table}

Các thao tác như sinh bảng ô chữ, kiểm tra đáp án, tính điểm, quản lý sao và lưu tiến độ là xử lý nội bộ của hệ thống sau khi người chơi thực hiện thao tác tương ứng. Vì vậy, các thao tác này không được mô hình hóa như một tác nhân ngoài trong biểu đồ use case.

\subsection{Biểu đồ use case tổng quát}
\label{subsection:2.2.1}
Ở mức tổng quát, hệ thống có các nhóm use case chính: quản lý tài khoản, chọn chủ đề, chọn màn chơi, giải ô chữ, sử dụng gợi ý, xem kết quả và theo dõi tiến độ. Trong nhóm quản lý tài khoản, người chơi có thể chọn tài khoản đã có hoặc tạo tài khoản mới. Trong nhóm chọn nội dung chơi, người chơi chọn chủ đề và màn. Trong nhóm giải ô chữ, người chơi chọn gợi ý, nhập đáp án, nhận phản hồi đúng sai và mở từ khi đúng. Trong nhóm kết quả và tiến độ, hệ thống hiển thị số sao đạt được, điểm, thời gian và trạng thái hoàn thành của màn. Hình~\ref{fig:usecase-tong-quat} mô tả các nhóm chức năng chính mà người chơi tương tác trong hệ thống.

\begin{figure}[H]
\centering
\resizebox{0.95\textwidth}{!}{\begin{tikzpicture}[
    font=\small,
    actorlabel/.style={align=center},
    usecase/.style={draw, ellipse, align=center, minimum width=4.0cm, minimum height=0.86cm, fill=white},
    boundary/.style={draw, rounded corners=4pt, minimum width=9.4cm, minimum height=7.5cm},
    assoc/.style={thick}
]
\coordinate (actor) at (0,0);
\draw[thick] (actor) ++(0,1.0) circle (0.23);
\draw[thick] (actor) ++(0,0.77) -- ++(0,-0.9);
\draw[thick] (actor) ++(-0.5,0.35) -- ++(1.0,0);
\draw[thick] (actor) ++(0,-0.13) -- ++(-0.45,-0.72);
\draw[thick] (actor) ++(0,-0.13) -- ++(0.45,-0.72);
\node[actorlabel] at ($(actor)+(0,-1.25)$) {Người chơi};

\node[boundary] (system) at (6.8,0) {};
\node[font=\bfseries, align=center] at ($(system.north)+(0,-0.36)$)
    {Trò chơi giải đố từ vựng dạng ô chữ};

\node[usecase] (account) at (6.8,2.45) {Quản lý tài khoản};
\node[usecase] (topic) at (6.8,1.62) {Chọn chủ đề};
\node[usecase] (level) at (6.8,0.79) {Chọn màn chơi};
\node[usecase] (play) at (6.8,-0.04) {Giải ô chữ};
\node[usecase] (hint) at (6.8,-0.87) {Sử dụng gợi ý};
\node[usecase] (result) at (6.8,-1.70) {Xem kết quả};
\node[usecase] (progress) at (6.8,-2.53) {Theo dõi tiến độ};

\coordinate (hub) at (2.0,0);
\draw[assoc] ($(actor)+(0.58,0.22)$) -- (hub);
\foreach \target in {account,topic,level,play,hint,result,progress}
    \draw[assoc] (hub) |- (\target.west);
\end{tikzpicture}
}
\caption{Biểu đồ use case tổng quát của hệ thống}
\label{fig:usecase-tong-quat}
\end{figure}

\begin{table}[H]
\centering
\begin{tabular}{|p{4cm}|p{9.8cm}|}
\hline
\textbf{Use case} & \textbf{Mô tả ngắn gọn} \\ \hline
Quản lý tài khoản & Người chơi chọn tài khoản cục bộ, tạo tài khoản mới hoặc xóa tài khoản không còn sử dụng. \\ \hline
Chọn chủ đề & Người chơi xem danh sách chủ đề, mỗi chủ đề hiển thị số màn đã hoàn thành trên tổng số màn. \\ \hline
Chọn màn chơi & Người chơi mở danh sách màn của chủ đề, xem số từ và trạng thái hoàn thành, sau đó bắt đầu màn. \\ \hline
Giải ô chữ & Người chơi chọn gợi ý, nhập đáp án, nhận phản hồi đúng sai và mở từ khi đáp án đúng. \\ \hline
Sử dụng gợi ý & Người chơi dùng sao để mở một chữ hoặc mở toàn bộ từ đang chọn. \\ \hline
Xem kết quả & Khi hoàn thành màn, người chơi xem điểm, số sao, thời gian, tỉ lệ đúng và nhiệm vụ đạt được. \\ \hline
Theo dõi tiến độ & Người chơi xem trạng thái hoàn thành ở danh sách chủ đề và danh sách màn chơi. \\ \hline
Ghi nhận tiến độ & Hệ thống tự động ghi nhận tiến độ sau khi người chơi hoàn thành màn hoặc có thay đổi kết quả cần lưu. \\ \hline
\end{tabular}
\caption{Danh sách use case tổng quát}
\label{table:usecase_overview}
\end{table}

\subsection{Biểu đồ use case phân rã}
\label{subsection:2.2.2}
Để làm rõ phạm vi chức năng, các use case tổng quát được phân thành bốn nhóm: tài khoản, chủ đề và màn chơi, gameplay, kết quả và tiến độ. Nhóm tài khoản phục vụ việc tách dữ liệu của nhiều người chơi trên cùng thiết bị. Nhóm chủ đề và màn chơi đưa người chơi từ màn hình chính tới màn chơi cụ thể. Nhóm gameplay bao gồm các thao tác lặp lại nhiều lần trong quá trình giải bảng ô chữ. Nhóm kết quả và tiến độ tập trung vào việc hiển thị kết quả, ghi nhận tiến độ và tải lại dữ liệu khi mở ứng dụng. Hình~\ref{fig:usecase-phan-ra} phân rã các use case theo bốn nhóm chức năng này. Các xử lý như kiểm tra đáp án, ghi nhận tiến độ và tải lại tiến độ là xử lý nội bộ được hệ thống thực hiện trong luồng tương tác.

\begin{figure}[H]
\centering
\resizebox{0.95\textwidth}{!}{\begin{tikzpicture}[
    font=\small,
    group/.style={
        draw,
        rounded corners=5pt,
        fill=gray!8,
        text width=4.7cm,
        minimum height=2.8cm,
        align=center,
        inner sep=8pt
    },
    actorlink/.style={thick, -{Stealth[length=2mm]}, gray!70}
]
\coordinate (actorbase) at (-3.2,0);
\draw[thick] (actorbase) ++(0,0.9) circle (0.22);
\draw[thick] (actorbase) ++(0,0.68) -- ++(0,-0.82);
\draw[thick] (actorbase) ++(-0.48,0.28) -- ++(0.96,0);
\draw[thick] (actorbase) ++(0,-0.14) -- ++(-0.42,-0.70);
\draw[thick] (actorbase) ++(0,-0.14) -- ++(0.42,-0.70);
\node[align=center] at ($(actorbase)+(0,-1.10)$) {Người chơi};

\node[group] (account) at (0.8,1.65) {
    \textbf{Tài khoản}\par\vspace{0.20cm}
    Tạo tài khoản\par
    Chọn tài khoản\par
    Xóa tài khoản
};
\node[group] (level) at (6.5,1.65) {
    \textbf{Chủ đề và màn chơi}\par\vspace{0.20cm}
    Chọn chủ đề\par
    Xem danh sách màn\par
    Bắt đầu màn chơi
};
\node[group] (gameplay) at (0.8,-1.85) {
    \textbf{Gameplay}\par\vspace{0.20cm}
    Chọn gợi ý\par
    Nhập đáp án\par
    Kiểm tra đáp án\par
    Dùng gợi ý
};
\node[group] (result) at (6.5,-1.85) {
    \textbf{Kết quả và tiến độ}\par\vspace{0.20cm}
    Hiển thị kết quả\par
    Ghi nhận tiến độ\par
    Tải lại tiến độ
};

\draw[actorlink] ($(actorbase)+(0.62,0.18)$) -- ($(account.west)+(-0.10,0)$);
\draw[actorlink] ($(actorbase)+(0.62,-0.18)$) -- ($(gameplay.west)+(-0.10,0)$);
\end{tikzpicture}
}
\caption{Biểu đồ use case phân rã theo nhóm chức năng}
\label{fig:usecase-phan-ra}
\end{figure}

\subsection{Quy trình chơi}
\label{subsection:2.2.3}
Quy trình nghiệp vụ chính của ứng dụng được tách thành hai phần để biểu đồ dễ đọc hơn. Phần thứ nhất mô tả luồng từ màn hình chính đến khi khởi tạo được một màn chơi hợp lệ. Người chơi chọn hoặc tạo tài khoản cục bộ, bấm Chơi, chọn chủ đề, sau đó chọn một màn trong danh sách màn của chủ đề đó. Hệ thống sinh bảng ô chữ trước khi chuyển sang màn chơi. Nếu không thể tạo màn hợp lệ, hệ thống hiển thị thông báo không thể tạo màn và quay lại danh sách màn. Nếu sinh bảng thành công, hệ thống hiển thị màn chơi và chuyển sang quy trình giải ô chữ. Hình~\ref{fig:activity-chon-man} mô tả quy trình từ chọn tài khoản đến khởi tạo màn.

\begin{figure}[H]
\centering
\resizebox{0.82\textwidth}{!}{\begin{tikzpicture}[
    font=\small,
    action/.style={draw, rounded corners=4pt, minimum width=4.5cm, minimum height=0.72cm, align=center, fill=white},
    decision/.style={draw, diamond, aspect=2.1, inner sep=1pt, minimum width=3.6cm, align=center, fill=yellow!10},
    startend/.style={draw, ellipse, minimum width=3.0cm, minimum height=0.68cm, align=center, fill=white},
    arrow/.style={-Stealth, thick},
    note/.style={font=\footnotesize}
]
\node[startend] (start) {Bắt đầu};
\node[action, below=0.45cm of start] (account) {Chọn hoặc tạo tài khoản};
\node[action, below=0.45cm of account] (play) {Bấm Chơi};
\node[action, below=0.45cm of play] (topic) {Chọn chủ đề};
\node[action, below=0.45cm of topic] (level) {Chọn màn chơi};
\node[action, below=0.45cm of level] (generate) {Hệ thống sinh bảng ô chữ};
\node[decision, below=0.75cm of generate] (ok) {Sinh bảng\\thành công?};
\node[action, right=2.2cm of ok] (fail) {Hiển thị thông báo\\không thể tạo màn};
\node[action, below=0.75cm of ok] (show) {Hiển thị màn chơi};
\node[action, below=0.45cm of show] (next) {Chuyển sang quy trình giải ô chữ};

\draw[arrow] (start) -- (account);
\draw[arrow] (account) -- (play);
\draw[arrow] (play) -- (topic);
\draw[arrow] (topic) -- (level);
\draw[arrow] (level) -- (generate);
\draw[arrow] (generate) -- (ok);
\draw[arrow] (ok) -- node[note, left] {Có} (show);
\draw[arrow] (show) -- (next);
\draw[arrow] (ok) -- node[note, above] {Không} (fail);
\draw[arrow] (fail.east) -- ++(0.6,0) |- (level.east);
\end{tikzpicture}
}
\caption{Activity diagram quy trình chọn màn và khởi tạo bảng}
\label{fig:activity-chon-man}
\end{figure}

Phần thứ hai mô tả quy trình giải ô chữ sau khi màn chơi đã được hiển thị. Người chơi chọn gợi ý, nhập đáp án và nhận phản hồi đúng sai. Nếu đáp án đúng, hệ thống mở từ, cộng điểm và cập nhật trạng thái màn. Nếu đáp án sai, hệ thống phản hồi để người chơi nhập lại. Khi gặp khó khăn, người chơi có thể dùng sao để mở một chữ hoặc mở cả từ nếu còn đủ sao. Khi tất cả các từ được hoàn thành, hệ thống tính điểm, tính sao nhiệm vụ, ghi nhận tiến độ và hiển thị kết quả. Hình~\ref{fig:activity-giai-o-chu} mô tả quy trình giải ô chữ, sử dụng gợi ý, hoàn thành màn và ghi nhận kết quả.

\begin{figure}[p]
\centering
\resizebox{!}{0.82\textheight}{\begin{tikzpicture}[
    font=\small,
    action/.style={draw, rounded corners=4pt, minimum width=4.2cm, minimum height=0.68cm, align=center, fill=white},
    decision/.style={draw, diamond, aspect=2.05, inner sep=1pt, minimum width=3.2cm, align=center, fill=yellow!10},
    startend/.style={draw, ellipse, minimum width=2.9cm, minimum height=0.64cm, align=center, fill=white},
    arrow/.style={-Stealth, thick},
    note/.style={font=\footnotesize}
]
\node[startend] (start) {Bắt đầu màn chơi};
\node[action, below=0.38cm of start] (select) {Người chơi chọn gợi ý};
\node[decision, below=0.58cm of select] (wantHint) {Người chơi\\dùng gợi ý?};
\node[decision, below=0.60cm of wantHint] (star) {Đủ sao?};
\node[action, right=1.9cm of star] (nostar) {Thông báo\\không đủ sao};
\node[action, below=0.55cm of star] (openhint) {Mở chữ hoặc từ,\\trừ sao};
\node[action, below=0.42cm of openhint] (guess) {Người chơi nhập đáp án};
\node[action, below=0.42cm of guess] (check) {Hệ thống kiểm tra đáp án};
\node[decision, below=0.60cm of check] (correct) {Đáp án\\đúng?};
\node[action, right=1.9cm of correct] (wrong) {Phản hồi sai\\và cho nhập lại};
\node[action, below=0.60cm of correct] (open) {Mở từ, cộng điểm,\\cập nhật trạng thái};
\node[decision, below=0.60cm of open] (done) {Tất cả từ\\hoàn thành?};
\node[action, below=0.60cm of done] (score) {Tính điểm,\\tính sao nhiệm vụ};
\node[action, below=0.38cm of score] (save) {Ghi nhận tiến độ};
\node[action, below=0.38cm of save] (result) {Hiển thị kết quả};
\node[startend, below=0.38cm of result] (end) {Kết thúc};

\draw[arrow] (start) -- (select);
\draw[arrow] (select) -- (wantHint);
\draw[arrow] (wantHint) -- node[note, left] {Có} (star);
\draw[arrow] (wantHint.west) -- ++(-0.75,0) node[note, above] {Không} |- (guess.west);
\draw[arrow] (star) -- node[note, left] {Có} (openhint);
\draw[arrow] (star) -- node[note, above] {Không} (nostar);
\draw[arrow] (nostar.east) -- ++(0.55,0) |- (select.east);
\draw[arrow] (openhint) -- (guess);
\draw[arrow] (guess) -- (check);
\draw[arrow] (check) -- (correct);
\draw[arrow] (correct) -- node[note, left] {Có} (open);
\draw[arrow] (correct) -- node[note, above] {Không} (wrong);
\draw[arrow] (wrong.east) -- ++(0.55,0) |- (guess.east);
\draw[arrow] (open) -- (done);
\draw[arrow] (done) -- node[note, left] {Có} (score);
\draw[arrow] (done.west) -- ++(-0.75,0) node[note, above] {Chưa} |- (select.west);
\draw[arrow] (score) -- (save);
\draw[arrow] (save) -- (result);
\draw[arrow] (result) -- (end);
\end{tikzpicture}
}
\caption{Activity diagram quy trình giải ô chữ và ghi nhận kết quả}
\label{fig:activity-giai-o-chu}
\end{figure}

\section{Đặc tả chức năng}
\label{section:2.3}
\subsection{Chọn chủ đề và màn chơi}
Use case này cho phép người chơi đi từ màn hình chính tới danh sách chủ đề và danh sách màn chơi. Tiền điều kiện là ứng dụng đã khởi động và dữ liệu chủ đề hợp lệ. Luồng chính gồm các bước: người chơi bấm nút chơi, hệ thống hiển thị danh sách chủ đề kèm tiến độ, người chơi chọn một chủ đề, hệ thống hiển thị các màn thuộc chủ đề đó, người chơi chọn màn và hệ thống khởi tạo trò chơi. Hậu điều kiện là màn chơi được hiển thị với bảng ô chữ, gợi ý và trạng thái ban đầu.

\begin{table}[H]
\centering
\begin{tabular}{|p{3.2cm}|p{10.8cm}|}
\hline
\textbf{Thuộc tính} & \textbf{Nội dung} \\ \hline
Tên use case & Chọn chủ đề và màn chơi \\ \hline
Tác nhân & Người chơi \\ \hline
Tiền điều kiện & Ứng dụng đã khởi động, dữ liệu chủ đề có ít nhất một chủ đề hợp lệ. \\ \hline
Luồng chính & Người chơi bấm Chơi; hệ thống hiển thị danh sách chủ đề; người chơi chọn chủ đề; hệ thống hiển thị danh sách màn; người chơi chọn màn; hệ thống sinh bảng và chuyển sang màn chơi. \\ \hline
Luồng phát sinh & Nếu không có chủ đề hợp lệ, hệ thống vô hiệu hóa nút bắt đầu và hiển thị trạng thái không có chủ đề. Nếu màn không sinh đủ số từ tối thiểu, hệ thống hiển thị thông báo không thể tạo màn và quay lại danh sách màn. \\ \hline
Hậu điều kiện & Màn chơi được khởi tạo, bảng ô chữ và danh sách gợi ý được hiển thị. \\ \hline
\end{tabular}
\caption{Đặc tả use case chọn chủ đề và màn chơi}
\label{table:uc_select_stage}
\end{table}

\subsection{Nhập đáp án}
Use case này cho phép người chơi giải một từ trong bảng ô chữ. Tiền điều kiện là màn chơi đang hoạt động và người chơi đã chọn một gợi ý. Người chơi nhập đáp án vào ô nhập liệu và bấm nút đoán. Hệ thống chuẩn hóa đáp án, so sánh với từ đích, sau đó phản hồi đúng hoặc sai. Nếu đúng, hệ thống mở toàn bộ chữ của từ, cộng điểm theo độ dài và thời gian giải, đồng thời cập nhật trạng thái hoàn thành của từ.

\begin{table}[H]
\centering
\begin{tabular}{|p{3.2cm}|p{10.8cm}|}
\hline
\textbf{Thuộc tính} & \textbf{Nội dung} \\ \hline
Tên use case & Nhập đáp án \\ \hline
Tác nhân & Người chơi \\ \hline
Tiền điều kiện & Người chơi đang ở màn chơi và đã chọn một gợi ý chưa hoàn thành. \\ \hline
Luồng chính & Người chơi nhập đáp án; hệ thống loại bỏ khoảng trắng thừa và chuẩn hóa chuỗi; hệ thống so sánh với đáp án đúng; nếu đúng, hệ thống mở từ, cộng điểm, cập nhật thống kê và lịch sử đoán. \\ \hline
Luồng phát sinh & Nếu chưa chọn gợi ý, hệ thống yêu cầu chọn gợi ý. Nếu ô nhập rỗng, hệ thống yêu cầu nhập từ đoán. Nếu đáp án sai, hệ thống tăng số lần đoán, ghi lịch sử và cho phép thử lại. \\ \hline
Hậu điều kiện & Từ được hoàn thành nếu đáp án đúng; trạng thái điểm, tiến độ và lịch sử đoán được cập nhật. \\ \hline
\end{tabular}
\caption{Đặc tả use case nhập đáp án}
\label{table:uc_guess}
\end{table}

\subsection{Sử dụng gợi ý}
Use case này hỗ trợ người chơi khi không tìm được đáp án. Tiền điều kiện là người chơi đã chọn một từ và còn đủ sao. Người chơi chọn mở một chữ hoặc mở cả từ. Hệ thống kiểm tra số sao, trừ sao nếu đủ, sau đó mở ký tự hoặc từ tương ứng. Nếu thao tác không tạo thay đổi, hệ thống hoàn lại sao. Hậu điều kiện là bảng ô chữ được cập nhật và số sao hiện tại thay đổi tương ứng.

\begin{table}[H]
\centering
\begin{tabular}{|p{3.2cm}|p{10.8cm}|}
\hline
\textbf{Thuộc tính} & \textbf{Nội dung} \\ \hline
Tên use case & Sử dụng gợi ý \\ \hline
Tác nhân & Người chơi \\ \hline
Tiền điều kiện & Người chơi đã chọn một gợi ý và tài khoản hiện tại còn đủ sao để dùng loại gợi ý tương ứng. \\ \hline
Luồng chính & Người chơi bấm mở một chữ hoặc mở cả từ; hệ thống kiểm tra số sao; hệ thống trừ sao; hệ thống mở chữ hoặc từ; giao diện cập nhật bảng, gợi ý và số sao. \\ \hline
Luồng phát sinh & Nếu người chơi chưa chọn gợi ý, hệ thống yêu cầu chọn gợi ý. Nếu không đủ sao, hệ thống hiển thị thông báo không đủ sao. Nếu từ đã được mở hết, hệ thống hoàn lại sao. \\ \hline
Hậu điều kiện & Một phần hoặc toàn bộ đáp án được mở, số sao và số lần dùng gợi ý được cập nhật. \\ \hline
\end{tabular}
\caption{Đặc tả use case sử dụng gợi ý}
\label{table:uc_hint}
\end{table}

\subsection{Hoàn thành màn và ghi nhận tiến độ}
Phần ghi nhận tiến độ được thực hiện như một xử lý nội bộ sau khi người chơi hoàn thành một màn, không phải là thao tác độc lập do một tác nhân hệ thống bên ngoài thực hiện. Khi toàn bộ từ trong màn đã được giải, hệ thống ghi nhận màn đã hoàn thành, số sao cao nhất, điểm cao nhất và thời gian tốt nhất vào tệp dữ liệu cục bộ. Dữ liệu được phân tách theo tài khoản để nhiều người chơi trên cùng máy có thể giữ tiến độ riêng.

\begin{table}[H]
\centering
\begin{tabular}{|p{3.2cm}|p{10.8cm}|}
\hline
\textbf{Thuộc tính} & \textbf{Nội dung} \\ \hline
Tên use case & Hoàn thành màn và ghi nhận tiến độ \\ \hline
Tác nhân kích hoạt & Người chơi, thông qua hành động hoàn thành màn. Hệ thống tự động thực hiện thao tác ghi dữ liệu sau khi màn chơi hoàn thành. \\ \hline
Tiền điều kiện & Người chơi hoàn thành toàn bộ từ của màn chơi hiện tại. \\ \hline
Luồng chính & Hệ thống tính số sao nhiệm vụ; cập nhật tổng sao; tăng số lượt chơi; ghi trạng thái hoàn thành của màn; lưu số sao cao nhất, điểm cao nhất và thời gian tốt nhất. \\ \hline
Luồng phát sinh & Nếu tệp lưu chưa tồn tại, hệ thống tạo cấu trúc dữ liệu mặc định. Nếu dữ liệu cũ không hợp lệ, hệ thống chuẩn hóa lại về cấu trúc hợp lệ. \\ \hline
Hậu điều kiện & Tệp tiến độ cục bộ được cập nhật theo tài khoản hiện tại. \\ \hline
\end{tabular}
\caption{Đặc tả xử lý hoàn thành màn và ghi nhận tiến độ}
\label{table:uc_save}
\end{table}

\section{Yêu cầu phi chức năng}
\label{section:2.4}
Ứng dụng cần có giao diện dễ sử dụng, các nút và thông tin trạng thái phải rõ ràng trên màn hình máy tính phổ thông. Bảng ô chữ cần tự điều chỉnh kích thước ô để hạn chế tràn khỏi vùng hiển thị. Thời gian sinh màn cần đủ nhanh để người chơi không phải chờ lâu khi bắt đầu một màn. Mã nguồn cần được chia thành các thành phần có trách nhiệm rõ ràng như sinh bảng, quản lý chủ đề, quản lý trạng thái, điều khiển giao diện và lưu dữ liệu. Dữ liệu tiến độ cần được lưu cục bộ ổn định và có khả năng phục hồi về mặc định nếu tệp lưu không hợp lệ.

\begin{table}[H]
\centering
\begin{tabular}{|p{2.7cm}|p{6.3cm}|p{4.8cm}|}
\hline
\textbf{Yêu cầu} & \textbf{Mô tả} & \textbf{Tiêu chí đánh giá} \\ \hline
Dễ sử dụng & Người chơi mới cần hiểu được luồng chơi qua các nút chính: Chơi, chọn chủ đề, chọn màn, chọn gợi ý và đoán đáp án. & Được kiểm tra trong thử nghiệm giao diện bằng cách đi từ menu tới màn chơi mà không cần hướng dẫn dài. \\ \hline
Hiệu năng & Việc sinh bảng cho một màn không được gây thời gian chờ dài; dữ liệu chủ đề không nên được tối ưu toàn bộ trước khi người chơi chọn màn. & Thời gian sinh bảng được đo trong phần kiểm thử hiệu năng; không tự đặt ngưỡng trước khi có kết quả đo. \\ \hline
Khả năng hiển thị & Bảng ô chữ cần nằm trong vùng chơi, các ô không quá lớn gây tràn và không quá nhỏ tới mức khó đọc. & Được kiểm tra ở một số kích thước cửa sổ, bảo đảm không chồng chữ và không tràn vùng chơi. \\ \hline
Độ ổn định dữ liệu & Nếu tệp lưu thiếu trường hoặc sai cấu trúc, hệ thống cần phục hồi về giá trị mặc định thay vì dừng chương trình. & Không phát sinh lỗi khi đọc tệp lưu thiếu trường hoặc có cấu trúc không hợp lệ trong kiểm thử. \\ \hline
Dễ bảo trì & Mã nguồn cần phân tách giữa dữ liệu chủ đề, thuật toán sinh bảng, trạng thái chơi, điều khiển giao diện và lưu dữ liệu. & Mã nguồn được tách theo nhóm trách nhiệm, hạn chế để toàn bộ logic trong một script giao diện. \\ \hline
Mở rộng nội dung & Hệ thống nên cho phép bổ sung chủ đề mới bằng cách thêm danh sách từ và gợi ý tương ứng trong bộ cung cấp màn chơi. & Có thể thêm chủ đề mới bằng cách bổ sung dữ liệu từ vựng và gợi ý, không cần sửa thuật toán sinh bảng. \\ \hline
\end{tabular}
\caption{Yêu cầu phi chức năng}
\label{table:nonfunctional}
\end{table}

Chương này đã xác định các chức năng chính và các ràng buộc phi chức năng của trò chơi giải đố từ vựng dạng ô chữ. Đây là cơ sở để lựa chọn công nghệ và thiết kế kiến trúc trong các chương tiếp theo.

\end{document}

